\documentclass[10pt,french]{book}
\input philippe2013

\newcounter{exoc}
\newenvironment{exoc}[1]{%
  \refstepcounter{exoc}\textbf{Exercice \theexoc} :\hfill {\textbf{#1}}\par
  \medskip}%
{\medskip}

\begin{document}

\pagestyle{empty}

\begin{center}
\renewcommand\arraystretch{1.5}
\begin{tabularx}{\textwidth}{|>\centering m{2.5cm}|>\centering X|>{\centering\arraybackslash} m{2.5cm}|}
	\hline
		1\iere \bsc{s.t.m.g.} & Mercredi 9 mars \np{2014} & \textbf{Proportions Probabilités} \\
	\hline
		\multicolumn{3}{|c|}{\bsc{Contrôle de mathématiques}} \\
	\hline
        \multicolumn{1}{|r}{\bsc{Nom}:} & \multicolumn{2}{l|}{} \\
		\multicolumn{1}{|r}{Prénom:} & \multicolumn{2}{l|}{} \\
	\hline
        \multicolumn{3}{|l|}{\bfseries Note et observations :} \\[1cm]
    \hline
\end{tabularx}\bigskip
\renewcommand\arraystretch{1}

{\itshape
La qualité et la précision de la rédaction seront prises en compte dans l'appréciation des copies.\par
Le barème est indicatif.}
\end{center}

\begin{exoc}{1,5 + 3,5 = 5 points}
    Soient $A$ et $B$ deux événements. On rappelle la formule suivante :
    \[p(A) + p(B) = p(A \cup B) + p(A \cap B).\]
    Dans un groupe de $250$ personnes, $100$ sont des femmes, $30$ sont des femmes exerçant la profession de chef d'entreprise et $120$ des personnes du groupe sont des chefs d'entreprise.\par
    On choisit au hasard une personne du groupe.\par
    On note $F$ et $C$ les événements suivants :\par
    $F$ : << la personne est une femme >>.\par
    $C$ : << la personne est chef d'entreprise >>.
    \begin{enumerate}
        \item En détaillant les calculs, déterminer $p(F)$, $p(C)$ et $p\left(\overline F\right)$. Donner les résultats \textbf{sous forme décimale}.
        \item Pour chacun des événements suivants, donner une description par une phrase simple puis déterminer \textbf{sous forme décimale} leur probabilité en justifiant précisément :
            \[C \cap F \qq C \cup F \qq \overline C \cap F \]
    \end{enumerate}
\end{exoc}\[*\]

\begin{exoc}{1,5 + 1 + 1,5 + 1 = 5 points}
\begin{minipage}{0.7\linewidth}
    Grâce à un système de détecteur, une borne de péage automatique peut délivrer des tickets à deux hauteurs différentes, selon le véhicule détecté, afin que le conducteur ne soit pas obligé de sortir de sa voiture.\par
    S'il s'agit d'une voiture, d'une moto ou d'une camionnette, le ticket sort en bas.\par S'il s'agit d'un camion, le ticket sort en haut.\par
    La société d'autoroute a étudié que l'une de ces bornes étaient défectueuses : de temps en temps, le ticket ne sort pas à la bonne hauteur, de façon indépendante pour chaque véhicule. D'après l'étude, la probabilité que le conducteur ne soit pas obligé de sortir de la voiture est égale à $0,9$.
\end{minipage}\qquad
\begin{minipage}{0.25\linewidth}
    \includegraphics[width=0.95\linewidth]{peage.eps}
\end{minipage}

    \begin{enumerate}
        \item On observe le comportement de la borne lors du passage de trois véhicules.\par Expliquer \textbf{précisément} pourquoi on peut modéliser la situation avec un schéma de Bernoulli.
        \item On appelle $A$ l'événement : << aucun conducteur n'est obligé de descendre de son véhicule pour récupérer le ticket >>. Calculer $p(A)$.
        \item On appelle $B$ l'événement : << un seul conducteur exactement descend de son véhicule pour récupérer le ticket >>. Calculer $p(B)$.
        \item On appelle $C$ l'événement : << au moins un conducteur descend de son véhicule pour récupérer le ticket >>.\par
        La société d'autoroute doit changer la borne défectueuse lorsque $p(C) \geqslant 25\%$.\par
        La borne étudiée doit-elle être changée ? Justifier précisément la réponse.
    \end{enumerate}
\end{exoc}

\begin{exoc}{4 + 3 + 1 + 1 + 1 = 10 pts}
Une enquête sur la propreté de leur ville a été menée auprès de $500$ personnes réparties de la manière suivante :
\begin{itemize}
    \item $30\%$ des personnes interrogées ont moins de $35$ ans ;
    \item $40\%$ ont entre $35$ et $50$ ans.
\end{itemize}

À la question << êtes-vous content de la propreté de votre ville  ? >>,
    \begin{itemize}
        \item $300$ des personnes interrogées ont répondu oui ;
        \item $120$ personnes de moins de $35$ ans ont répondu oui ;
        \item $70\%$ des personnes de plus de $50$ ans ont répondu non.
    \end{itemize}

\begin{enumerate}
    \item À l'aide des informations données, compléter le tableau suivant :
    \begin{center}
    \renewcommand\arraystretch{2.5}
        \begin{tabular}{|c|c|c|c|>{\centering\arraybackslash}p{2cm}|}
            \hline
                & moins de $35$ ans & entre $35$ et $50$ ans & plus de $50$ ans & Total \\
            \hline
                Réponse Oui & & & & \\
            \hline
                Réponse Non & & && \\
            \hline
                Total & & & & \\
            \hline
        \end{tabular}
    \renewcommand\arraystretch{1}
    \end{center}

    \item Donner sous forme d'un pourcentage :
        \begin{enumerate}
            \item la proportion $p_1$ de personnes qui ont répondu Non ;
            \item la proportion $p_2$ de personnes qui ont entre $35$ et $50$ ans \textbf{et} qui ont répondu Non ;
            \item  la proportion $p_3$ de personnes qui ont \textbf{moins} de $50$ ans \textbf{et} qui ont répondu Oui.
        \end{enumerate}
    \item Quelle est la proportion $p_4$ de personnes ayant moins de $35$ ans parmi toutes les personnes qui ont répondu Oui ?
    Donner le résultat sous forme d'un pourcentage.
    \item On appelle $N$ la population des personnes qui ont répondu Non et $T$ la population des personnes âgées entre $35$ et $50$ ans.\par
    En détaillant la démarche, calculer $p_{(T \cup N)}$. Donner le résultat en pourcentage.
    \item On rappelle que $30\%$ des personnes interrogées ont moins de 35 ans. Parmi ces personnes, $48\%$ sont des femmes.\par
    Déterminer le pourcentage de femmes de moins de 35 ans sur l'ensemble des personnes interrogées.
\end{enumerate}
\end{exoc}

\end{document} 